\section{Main Part}\label{Sec:MainPart}

\subsection{Selected design figures from the thesis}
The figures adapt the thesis schematics with English labels and generalised component names.
The opening Figure~\ref{fig:paper-orchestration}, placed before the abstract, compares the three orchestration variants.
Figure~\ref{fig:paper-architecture} summarises the integration architecture,
and Figure~\ref{fig:paper-dialogue-state} shows how state is carried across dialogue turns.
The architecture spans both columns to retain readable labels;
the dialogue-state diagram is arranged vertically within one column.

\begin{figure*}[t]
    \centering
    \resizebox{\textwidth}{!}{% Adapted from the author's thesis, 04_Systementwurf_Architektur_neu.tex,
% fig:servicearchitektur_mcp. English labels and generalised service roles;
% internal identifiers, data-system names and authorisation services omitted.
\begin{tikzpicture}[
  font=\fontsize{8.5}{10}\selectfont, >=Latex,
  box/.style={draw,rounded corners=2pt,align=center,inner sep=5pt,
    minimum height=1.05cm,text width=2.8cm,fill=white},
  connection/.style={<->,thick},
  every node/.style={outer sep=1pt}
]
  \draw[rounded corners=3pt,draw=black!55,fill=blue!3]
    (1.55,-1.25) rectangle (10.55,3.25);
  \node[anchor=north west,font=\fontsize{9}{11}\selectfont\bfseries]
    at (1.7,3.12) {MCP-based extension};
  \node[box,text width=1.2cm,fill=blue!10] (user) at (0,0) {User};
  \node[box,text width=3cm] (client) at (3.65,0)
    {\textbf{Dialogue client / host}\\Dialogue state\\MCP client};
  \node[box,text width=2.3cm,minimum height=0.75cm] (llm) at (3.65,2)
    {\textbf{Language model}};
  \node[box,text width=2.6cm] (server) at (8.6,0)
    {\textbf{MCP server}\\Defined tools\\Validation};
  \node[box,text width=2.9cm] (services) at (13.4,0)
    {\textbf{Existing service APIs}\\Scheduling and data services};
  \draw[connection] (user.east) -- node[above,font=\fontsize{8}{9}\selectfont] {Dialogue} (client.west);
  \draw[connection] (client.north) -- node[right,align=left]
    {Interpretation and\\state updates} (llm.south);
  \draw[connection] (client.east) -- node[above] {MCP} (server.west);
  \draw[connection] (server.east) -- node[above] {API calls} (services.west);
\end{tikzpicture}
}
    \caption{Generalised architecture of a scheduling assistant using the Model Context Protocol (MCP).
    The user interacts with a dialogue client that maintains state and asks a language model to interpret inputs.
    The client invokes tools on the MCP server, which validates requests and accesses existing scheduling and data services through their application programming interfaces (APIs).
    Double-headed arrows represent requests and responses, not unrestricted model access to services.
    Internal system names and deployment details are omitted.}
    \label{fig:paper-architecture}
\end{figure*}

\begin{figure}[tb]
    \centering
    \resizebox{\columnwidth}{!}{% Adapted from the author's thesis, 05_Realisierung_neu.tex,
% fig:dialogverarbeitung_workcontext. Reflowed for one paper column.
% The implementation-specific context name is replaced by dialogue state.
\begin{tikzpicture}[
  font=\fontsize{8.5}{10}\selectfont, >=Latex,
  box/.style={draw,rounded corners=2pt,align=center,inner sep=4pt,
    minimum height=0.8cm,text width=2.5cm,fill=white},
  arrow/.style={->,thick},
  every node/.style={outer sep=1pt}
]
  \node[box,fill=blue!10] (message) at (0,0) {Current user input};
  \node[box] (state) at (3.9,0) {Previous\\dialogue state};
  \node[box,text width=3.5cm,minimum height=1cm] (process) at (1.95,-1.6)
    {Dialogue processing\\according to the selected\\orchestration variant};
  \node[box] (response) at (0,-3.2) {System response};
  \node[box] (updated) at (3.9,-3.2) {Updated\\dialogue state};
  \draw[arrow] (message.south) -- (process.north west);
  \draw[arrow] (state.south) -- (process.north east);
  \draw[arrow] (process.south west) -- (response.north);
  \draw[arrow] (process.south east) -- (updated.north);
  \draw[arrow] (updated.east) -- (5.7,-3.2) --
    node[right,rotate=90,anchor=south] {Next dialogue turn} (5.7,0) -- (state.east);
\end{tikzpicture}
}
    \caption{State handling across successive dialogue turns.
    The processing stage combines the current user input with information retained from previous turns, following the selected orchestration variant.
    It produces a response for the user and an updated dialogue state.
    The right-hand feedback arrow carries that state into the next turn, so previously established information can be reused.}
    \label{fig:paper-dialogue-state}
\end{figure}

\subsection{Testing some math}

A regular formula

\begin{equation}
    x^2 + y^2 = z^2 
\end{equation}
and another one
\begin{equation}
    \left\Vert \vec{x} \right\Vert = \sqrt{\sum_{i=0}^{n}{x_i^2 + y_i^2}}.
\end{equation}

Here is something with matrices:
\begin{equation}
\bf{A}=
\begin{bmatrix}
    a & b & c\\
    d & e & f\\
\end{bmatrix}^\top.
\end{equation}

An some inline stuff $\sum a_i=\bf{B}$.

\begin{table}
    \caption[List of Symbols]{Keep a list of symbols in your paper.}
    \label{tab:ListOfSymbols}
    \resizebox{\columnwidth}{!}
    {
    \begin{tblr}{ll}
    \toprule
    \SetCell[c=2]{c}Variables used in this paper\\
    \midrule
    \bf{Symbol}               & \bf{Meaning}\\
    \midrule
    $\mat{M}$                 & model-view matrix \\
    $\mat{P}$                 & projection matrix \\
    $\nvec{n}$                & a normal vector \\
    $V$                       & number of vertices \\    
    $T$                       & number of triangles \\    
    $\vec{l}$                 & light direction vector \\
    $\pnt{x}$                 & point to be lit \\
    \bottomrule
    \toprule
    \SetCell[c=2]{c}{Types used in this paper}\\
    \midrule
    \bf{Notation}               & \bf{Meaning}\\
    \midrule
    $a$                         & scalar\\
    $N$                         & scalar constant\\
    $\alpha$                    & angle\\
    $\vec{x}$                   & column vector\\
    ${x}_i$                     & vector component\\
    $\nvec{x}$                  & normalized vector \\
    $\vec{x}=
    \begin{bmatrix}
        {x}_0\\
        {x}_1\\
        \vdots\\
        {x}_{n-1}
    \end{bmatrix}$              & column vector\\
    $\mat{A}=
    \begin{bmatrix}
        a_{0,0}   & a_{0,1}   & \dots  &  a_{0,m-1}   \\
        a_{1,0}   & a_{1,1}   & \dots  &  a_{1,m-1}   \\
        \vdots    & \vdots    & \ddots &  \vdots      \\
        a_{n-1,0} & a_{n-1,1} & \dots  &  a_{n-1,m-1} \\
    \end{bmatrix}
    $                          & matrix \\
    $\vec{a}_{:,j}$             & $j$-th column of matrix $\mat{A}$\\
    ${a}_{i,j}$                 & element at row $i$, column $j$\\
    $\quat{q}$                  & quaternion \\
    $\pnt{p}$                   & point \\
    $s_i$                       & $i$-th element of a series of scalars\\
    $\vec{x}^{\left[i\right]}$  & $i$-th element of a series of vectors\\
    $\mat{A}^{\left[i\right]}$  & $i$-th element of a series of matrices\\
    $\mathcal{A}$               & set\\
    \bottomrule                                    
    \toprule
    \SetCell[c=2]{c}{Operators used in this paper}\\
    \midrule
    \bf{Operation}       & \bf{Meaning}\\
    \midrule
    $\vec{a}^\top$           & transpose, row vector \\
    $\vec{a}^\top\vec{b}$    & inner product \\
    $\vec{a} \times \vec{b}$ & cross product \\    
    \bottomrule                                    
\end{tblr}

    }
\end{table}
